\section{Methods}

\subsection{Model Overview}

Since the prediction target is a binary success/failure for given drug-disease indication, we
 evaluated several supervised models for binary classification, 
including Logistic Regression, 
k-Nearest Neighbors (kNN), 
Support Vector Machines (SVM), 
Random Forest, 
Gradient Boosting, 
and XGBoost.

Logistic regression estimates probability of success using sigmoid-transformed
 linear combinations of the input features. Its biggest advantage is high interoperability
 since feature waits indicate their importance, however it is limited to linearly
 separable problems.

k-Nearest Neighbors is instance based model, making predictions by looking at labels
of the closest points in the feature space. Therefore, there is no actual training phase and the assumption behind this model is that similar observations
tend to have the same label. Nevertheless it is very sensitive to noise and the curse of dimensionality, 
as well as poorly interpretable (only by samples that voted).

Support Vector Machines are fitting a decision boundary, maximizing the margin between classes.
Although it may be similar to regularized linear regression, the usage of kernel functions
allows fitting of non-linear decision boundaries.

Tree-based ensemble learning models leverage hierarchical decision rules
 and combine prediction of multiple models. A random forest aggregates 
 multiple decision trees based on random subsets of variables and observations, 
 this decreases the overfitting so characteristic of decision trees. Gradient boosting and
 XGBoost instead build trees sequentially, where each new tree learns errors of the previous ones.
 XGboost is an optimized version of Gradient Boosting, but has more parameters to tune.

As for graph-based approach, we chose to use a Heterogenous Graph Neural Network (GNN) link classifier from PyTorch Geometric module.
Instead of learning patterns based on tabular features only, it learns representations by propagating information
across connected nodes and edges, enabling the model to capture more complex relationships.
This approach avoids manual feature engineering, but requires a large number of samples and 
rich informative connections in a graph, and, as other Artificial Neural Networks, prone to overfitting.

\subsection{Implementation Details}

All numeric features are centered and scaled. No imputation was needed since 
observations with missing values corresponding to biopolymer drugs were excluded.

Most of the mentioned models are imported from \verb|sklearn| module, except for XGBoost
imported from \verb|xgboost| and GNN from \verb|torch_geometric| modules. In order to perform
hyperparameter tuning, we used a grid-search cross validation \verb|GridSearchCV| with search spaces indicated 
in Table~\ref{tab:hyperparameters}.

Since most of the clinical trials and thus corresponding indications usually fail,
the classes are imbalanced with failures being several times more abundant (around 3 times for our data).
So we used options for class balancing where it was possible (all models except for kNN and Gradient Boosting).

\begin{table}[H]
    \centering
    \footnotesize
    \caption{Hyperparameter search spaces for each classifier.}
    \label{tab:hyperparameters}
    \rowcolors{2}{gray!15}{white}
    \begin{tabular}{p{4cm} p{10cm}}
        \toprule
        \textbf{Model} & \textbf{Hyperparameters} \\
        \midrule
        k-Nearest Neighbors & \texttt{n\_neighbors} $\in \{5, 10, 20, 30\}$, \texttt{weights} $\in \{\text{uniform}, \text{distance}\}$, \texttt{p} $\in \{1, 2\}$ \\
        Logistic Regression & \texttt{C} $\in \{0.01, 0.1, 1, 10, 100\}$, \texttt{l1\_ratio} $\in \{0, 0.25, 0.5, 0.75, 1\}$ \\
        SVM (linear) & \texttt{C} $\in \{0.01, 0.1, 1, 10\}$ \\
        SVM (polynomial) & \texttt{C} $\in \{0.1, 1, 10\}$, \texttt{degree} $\in \{2, 3, 4\}$, \texttt{gamma} $\in \{\text{scale}, \text{auto}\}$ \\
        SVM (RBF) & \texttt{C} $\in \{0.1, 1, 10\}$, \texttt{gamma} $\in \{\text{scale}, \text{auto}\}$ \\
        Random Forest & \texttt{n\_estimators} $\in \{100, 300, 500\}$, \texttt{max\_depth} $\in \{\text{None}, 10, 20\}$, \texttt{min\_samples\_split} $\in \{2, 5\}$, \texttt{min\_samples\_leaf} $\in \{1, 2\}$ \\
        Gradient Boosting & \texttt{n\_estimators} $\in \{100, 300\}$, \texttt{max\_depth} $\in \{3, 5\}$, \texttt{learning\_rate} $\in \{0.01, 0.05, 0.1\}$, \texttt{subsample} $\in \{0.8, 1.0\}$ \\
        XGBoost & \texttt{n\_estimators} $\in \{100, 300\}$, \texttt{max\_depth} $\in \{3, 5\}$, \texttt{learning\_rate} $\in \{0.01, 0.05, 0.1\}$, \texttt{subsample} $\in \{0.8, 1.0\}$, \texttt{colsample\_bytree} $\in \{0.8, 1.0\}$, \texttt{gamma} $\in \{0, 0.1\}$ \\
        \bottomrule
    \end{tabular}
\end{table}


\subsection{Evaluation Strategy}

Drug-disease indication success prediction requires avoiding temporal data leak - 
in order to simulate fairly the real application, we need to use older research results
to predict success of new indications. Thus, we performed a train-test split based on
on start date of the first clinical trial for each indication. Taking 15\% for test set 
results in temporal splitting by August 2014.

The training set is further used for hyperparameter tuning using 
grid search with 5-fold cross-validation and model selection based 
on cross-validation performances. 

Since we deal with imbalanced classes, we must not use accuracy as evaluation metric since
it will underemphasize incorrect prediction of minority (success) class. We thus used F1 score
(harmonic mean of precision and recall) for model hyperparameter tuning, selection and evaluation, 
as well as precision-recall curves instead of ROC curves for analyzing how clarification result varies 
based on decision threshold (it works only for models yielding class probabilities).

